% =============================================================================
% 10. CONCLUSION
% Target: 0.5 page (~250 words)
% =============================================================================

\section{Conclusion}
\label{sec:conclusion}

% ---------- ¶1 Recap ----------
We introduced \textsc{PATHCAST}, a formally specified pipeline for
process-aware stochastic forecasting of the software development
lifecycle. The pipeline integrates event log construction, process
mining, absorbing Markov chain estimation, process-aware Monte Carlo
simulation, and a machine-learning enhancement component, under
contract-based stage composition (Table~\ref{tab:contracts}). The
contribution is not the invention of new individual techniques but
their formal integration into an SDLC-specific pipeline with provable
properties: pipeline correctness under stated preconditions
(Theorem~\ref{thm:pipeline-correctness}), asymptotic statistical
consistency (Theorem~\ref{thm:consistency}), end-to-end distributional
convergence with vanishing calibration error
(Theorem~\ref{thm:pipeline-convergence}), and an information-theoretic
advantage of process awareness over throughput-based Monte Carlo
(Theorem~\ref{thm:information}). A closed-form numerical example on a
six-state SDLC chain (Section~\ref{sec:numerical-example})
cross-validates the analytical and Monte Carlo computations to within
$1\%$ on every quantity for which a closed form exists.

% ---------- ¶2 Limitations ----------
The theoretical results rest on five explicit assumptions
(Table~\ref{tab:assumptions}): first-order Markov dynamics,
stationarity, case independence, sojourn--history independence, and
complete case observation. Each assumption is reported together with
the consequence of its violation and the mitigation adopted in the
pipeline (state augmentation, sliding-window estimation, drift
detection, or inclusion of right-censored cases). None of the
assumptions is innocuous in practice; their empirical impact on
forecast quality is investigated in the companion paper.

% ---------- ¶3 Future work ----------
Three directions are anticipated. \emph{First,} the empirical
validation of \textsc{PATHCAST} across 48 open-source repositories
with comparison against four baselines and statistical analysis
(Section~\ref{sec:evaluation-plan}) is reported in the companion
paper. \emph{Second,} the relaxation of the first-order Markov
assumption via state augmentation or history-conditioned semi-Markov
models is a natural theoretical extension, motivated by the rework
patterns that the present model captures only indirectly.
\emph{Third,} integrating resource and queueing effects (relaxing the
case-independence assumption) would extend the pipeline from
case-level to team-level forecasting and connect this line of work to
the classical queueing theory used in software process performance
analysis.
